\section{Theory requirements and numerical controls}\label{sec:tt-theory}
A microscopic prediction must connect a specified many-nucleon Hamiltonian and treatment of cluster antisymmetry to elastic-channel $S$-matrix elements. Cluster potential models offer useful controlled parameterizations but do not in general yield a unique underlying six-body interaction.

For multichannel dynamics, separate the elastic amplitude from contributions of open or virtual rearrangement channels and use physically appropriate asymptotic Coulomb conditions. The logical structure of three-body scattering equations is exemplified by the foundational work of \textcite{Faddeev1961}; its citation here is methodological context, \emph{not} a claim that the original work calculated $t+t$ scattering.

A reproducible numerical comparison should report: (i) angular-momentum and model-space convergence, (ii) Coulomb asymptotic treatment, (iii) sensitivity to interaction and regulator choices, (iv) threshold placement, and (v) a mapping from theoretical amplitudes to the \emph{measured} observable including instrumental response.

\subsection{Connection to correlation measurements}
A correlation function may provide an additional constraint only with an independently justified source model and treatment of quantum statistics and final-state interactions. For the conceptual framework see \textcite{Koonin1977,Lednicky1982}. These are foundational references, not direct $t+t$ data papers.
